\documentclass[11pt]{article}
\usepackage{amsmath, amssymb}
\usepackage[margin=1in]{geometry}
\usepackage{enumitem}

\title{Projection Tactics: Experiment Descriptions}
\author{}
\date{}

\begin{document}
\maketitle

\noindent All experiments use contrastive training pairs $(x_i, y_i^+, y_i^-)$ and extract response-token activations $a^l$ at layer $l$.

% ------------------------------------------------------------------
\section{DiffMean Projection Analysis}
\label{sec:diffmean}

Projects test-time activations onto the DiffMean steering direction and checks correlation with unsteered behavior scores.

\begin{align}
    \mathbf{v}^l &= \boldsymbol{\mu}_+ - \boldsymbol{\mu}_-, \quad \hat{\mathbf{v}} = \mathbf{v}/\|\mathbf{v}\| \\
    c_+ &= \boldsymbol{\mu}_+ \cdot \hat{\mathbf{v}}, \quad c_- = \boldsymbol{\mu}_- \cdot \hat{\mathbf{v}} \\
    m &= \tfrac{c_+ + c_-}{2}, \quad h = \tfrac{c_+ - c_-}{2} \\
    \tilde{p}_j &= \frac{\bar{a}^l(x_j, y_j) \cdot \hat{\mathbf{v}} - m}{h}
\end{align}

Under this normalization $c_+ \to +1$, $c_- \to -1$. Tests whether $\rho(\tilde{p},\, \text{score}) \gg 0$.

% ------------------------------------------------------------------
\section{LDA Projection Analysis}
\label{sec:lda}

Same projection and centroid normalization as Section~\ref{sec:diffmean}, but replaces the DiffMean direction with the Fisher LDA direction:
\begin{equation}
    \mathbf{w}_\text{LDA} = \boldsymbol{\Sigma}_W^{-1}(\boldsymbol{\mu}_+ - \boldsymbol{\mu}_-)
\end{equation}
where $\boldsymbol{\Sigma}_W$ is the pooled within-class covariance. Accounts for the shape of activation clusters, not just their means.

% ------------------------------------------------------------------
\section{Logistic Probe Projection Analysis}
\label{sec:probe}

Same framework as Section~\ref{sec:diffmean}, but projects onto the weight vector of a logistic regression probe:
\begin{equation}
    \mathbf{w}_\text{probe} = \arg\min_{\mathbf{w}} \sum_i \mathcal{L}_\text{BCE}\!\bigl(\sigma(\mathbf{w} \cdot \bar{a}^l_i + b),\; y_i\bigr)
\end{equation}
where positive activations are labeled 1, negative labeled 0. Optimizes a classification objective over all training activations.

% ------------------------------------------------------------------
\section{Directional Agreement (Steering Vector)}
\label{sec:dir-agree}

Measures how well each prompt's activation difference aligns with the steering vector.
\begin{align}
    \boldsymbol{\Delta}_i &= \bar{a}^l(x_i, y_i^+) - \bar{a}^l(x_i, y_i^-) \\
    \mathbf{s}^l &= \frac{1}{N}\sum_i \boldsymbol{\Delta}_i \\
    \text{score}_i &= \cos(\boldsymbol{\Delta}_i,\; \mathbf{s}^l)
\end{align}
A narrow, high-mean distribution indicates the steering vector consistently captures individual examples' behavioral axes.

% ------------------------------------------------------------------
\section{Directional Agreement (Pairwise)}
\label{sec:pairwise}

Instead of comparing each $\boldsymbol{\Delta}_i$ to the steering vector, compares all pairs of per-prompt deltas:
\begin{equation}
    \text{score}_{ij} = \cos(\boldsymbol{\Delta}_i,\; \boldsymbol{\Delta}_j), \quad i < j
\end{equation}
High pairwise agreement means individual prompts shift activations in a consistent direction, independent of the steering vector definition.

% ------------------------------------------------------------------
\section{Token-Position Directional Agreement}
\label{sec:token-pos}

Per-token version of Section~\ref{sec:dir-agree}. Instead of averaging activations first, computes a difference vector at each token position and averages the cosine similarities.
\begin{align}
    \delta_i^t &= a^l_+(x_i, t) - a^l_-(x_i, t), \quad t = 1, \dots, \min(T_+, T_-) \\
    \mathbf{s}^l &= \frac{1}{N}\sum_i \bigl[\bar{a}^l_+(x_i) - \bar{a}^l_-(x_i)\bigr] \\
    \text{score}_i &= \frac{1}{T_{\min}} \sum_t \cos(\delta_i^t,\; \mathbf{s}^l)
\end{align}
The steering vector $\mathbf{s}^l$ uses all tokens (no truncation); the per-token cosine similarities are computed only up to $\min(T_+, T_-)$.

% ------------------------------------------------------------------
\section{Activation Spread}
\label{sec:spread}

Measures how dispersed a single prompt's response-token activations are, computed separately for positive and negative responses.
\begin{equation}
    \text{spread}_i = \frac{1}{\binom{T}{2}} \sum_{t < t'} \bigl[1 - \cos(a^l_t,\; a^l_{t'})\bigr]
\end{equation}
Higher spread means token-position activations point in more diverse directions. Reported per-prompt per-layer, and averaged across prompts.

\end{document}
